\section{把中枢放回可走的路上}

\subsection{魏：邺、洛阳与关中的不同时钟}

洛阳的宫门使魏的朝廷看起来很集中，但一项北方政策往往先在邺、长安和州郡的
仓中有了自己的速度。魏文帝选择洛阳，是要把受禅后的尚书台、宿卫和旧汉官僚
接在一处；《三国志·文帝纪》所能确认的，是都城和诏令的次序。它不能替我们
看见每一个被留任的官吏为何接受新印信。对近畿以外的人，朝廷是否已经换名，
要等到任命、征发和诉讼的文书沿道路抵达，才成为一个必须回答的日常问题。

这条差别在曹叡朝的关中尤其清楚。蜀军从汉中向陇右伸出时，魏廷既要把援军
送过关中，也要避免为了西线而抽空淮南。张郃等将领的行动被《三国志》写进
胜负年序，真正使援军可以行动的却是关隘、驿路、马匹和能够暂借粮食的地方
节点。朝廷可以命令一支军队向西，不能命令春雨、栈道和运夫同时服从。于是
魏的优势不是没有距离，而是在距离已经造成麻烦时，仍比蜀汉更可能从若干方向
调来替手。这个优势也有价格：边地的州郡和运输者承担的劳役，很少像将领的
姓名一样进入本纪。

二三〇年的伐蜀准备被雨水和道路拦下，正好抵抗一种过于平滑的北方叙事。若
只从洛阳读诏书，行动似乎已经开始；若把山道也放进句子，便会看见军令抵达
关中后还要等待车、畜和可通行的栈道。魏并不是因为“不愿战”而停下，也不能
由一次停下推论其军力衰弱。对皇帝和中枢官员来说，取消行动是承认资源不能
在预定时日合拢；对沿线承担准备的人来说，已经发生的征发未必随着撤令完全
消失。战争的后方因此不是战场之外的空白，而是决定何时可以不打的一组成本。

\subsection{品第与名门：一套入口，不是一张社会总表}

陈群所议九品官人法常被写成魏晋门阀的开端。把它放回黄初初年的洛阳，较能
看见它面对的实际问题：新朝需要把察举、乡里声望和中枢任用接成可操作的次序，
而不是每逢官缺都重新从战争中的流散人群辨认资格。《三国志·陈群传》保留了
制度名称和陈群的政治位置；它不能给出各州品评的完整名单，更不能证明每一位
被评为上品的人都以同一种家世进入官署。

品评的力量在于它让地方评价可以走向中央，局限也正在这里。谁能被看见、谁
能由乡里大姓或州郡官员作保、谁的迁徙使旧有关系断裂，都会改变一张名册的
形状。战乱后从河北、关中或荆州进入新政权的人，并非天然处在同一条起跑线上。
制度给洛阳一个筛选入口，也把地方社会已有的不均衡带到入口之前。因而不能把
九品简单写成魏“发明贵族”，也不应把它说成完全中性的行政技术；它是在有限
资讯和需要合作的地方精英之间作出的选择。

正始石经的残片提供另一种中枢可见性。石面上的异体字和经文校读，让人看见
太学并非只有抽象的儒学声望，而有刻石、抄写、阅读和观看的场所。残石不能
告诉我们远方县学的课程，更不能证明农户如何读经；它仍提醒我们，国家也在
用可公开校对的文本训练官员的共同语言。品第处理“谁可入仕”，石经处理
“何种文字可被援引”；两者都在洛阳运作，却都要经由比宫门更长的地方关系。

\subsection{蜀：一座盆地怎样供应一条北线}

成都的优势从来不是一块自动可用的粮仓。都江堰灌溉与成都平原的富庶是后世
理解盆地的重要线索，但三国材料并没有留下可让我们逐年计算收成、赋额和军粮
的总账。比较稳妥的说法是：成都能够支撑长期官署和军队，正因为水利、聚落、
地方官与转输网络在这里比山口以外更容易接合；一旦粮食要越过汉中，优势便
被路程重新分割。

诸葛亮驻汉中并非把一个“前进基地”放到地图上便告完成。汉中城、褒斜道和
沿线仓口必须持续有人守护，来自盆地的粮食也必须在山谷、关隘与前线之间反复
转手。街亭失守后撤军，常被归成马谡的个人成败；可是守一处山口的决定所以
沉重，恰在于蜀军没有另一条无成本的补给线。将领的判断会改变道路是否还能
使用，道路的脆弱又会把判断的后果放大。这里不需要虚构帐篷里的争论，也能
理解为什么一次据点失守足以迫使远处的成都接受撤军。

蒋琬、费祎时期的谨慎同样不宜写成没有行动。维持汉中守备、安置退军、与吴
保持交往、让成都的任官与供应不被一次北伐耗尽，都是选择。它们不像出祁山那
样留下醒目的战场名称，却解释了蜀汉为何能在诸葛亮死后继续存在。二五三年
费祎被刺，损失的不仅是一位名臣，而是一个能在北向主张、宫廷日常与有限人力
之间调节的人。其后姜维更频密地北向，不应被倒写为费祎若在必能阻止；较可见
的是，成都用于限制与协调的一个入口消失了。

\subsection{南中不是“后方”这个词可以收纳的地方}

《华阳国志》与《三国志·蜀书》使我们知道诸葛亮二二五年南下、郡县关系被
重整，也保存了许多后来被反复讲述的人物和故事。它们的叙事方向首先从成都
出发：南中被写成蜀汉需要重新连上的边缘。若以此便断言当地社会已经被完全
驯服，反而会把材料的立场当作事实。更适合追问的是，成都为什么需要那里：
南方的郡治、道路、金属与人力如何影响盆地可以抽调多少资源，以及地方首领
为何能成为传令、征发和安置不可绕过的中介。

“平定”在朝廷文字中可以很快结束，在山地和聚落之间却只能是不断重作的关系。
州郡官员需要翻译、向导、仓储与当地承认；地方群体也会评估新来的军队带来
的是保护、征调还是风险。没有连续的地方文书，不能替他们写出统一的意见。
保留这种沉默不是削弱蜀汉的叙事，而是让读者明白北伐所依赖的后方本身也有
自己的政治时间。
